\chapter{Tests}
\ChapterLabel{Tests}
\section{Chapter Responsibility}
\label{Section:TestsResponsibility}
Define repeatable procedures that produce objective evidence.

\noindent\textbf{Expected result:} Executable test specifications with observable pass criteria.

\section{How to Use This Chapter}
Define repeatable test cases, levels, environments, data, entry conditions, procedures, expected results, and result recording. Tests produce evidence; they do not replace verification or validation.
\section{Week 2 Planned Checks}
\label{Section:WeekTwoTests}
These are candidate test specifications. Jack Marino configures and records
the hardware, host, driver, versions, operating envelope, independent reference,
instrument calibration/uncertainty, and applicable parameters from
\hyperref[Table:RequirementParameters]{the decision register} before running
each check. Resolve the mapped requirement's approval dependencies first.
For each run retain inputs, raw logic/timing traces or records, calculations,
observed results, reviewer, and disposition. A missing bound or reference
blocks a pass/fail judgment; no check below has been executed.
Main and supporting scope follow \hyperref[Table:MainBaseline]{the selected requirements baseline}. All checks remain Pending; evidence is planned.
\paragraph{Permitted planned evidence environments.} Controlled replay or synthetic records may exercise local decoding, ordering, recording/replay and software fault/status logic (for example TC-016, TC-050 and TC-051). Such observations cover only the declared software boundary. Physical ready-to-publication latency, ADC validity, source timestamp accuracy, disconnect recovery and the approved OS/driver matrix need the specified target and independent reference. TC-032 and TC-054 may use a simulated mission consumer only after its contract, schema/profile, decoder and controlled environment are approved; an OMS conformance disposition needs that approved integration evidence. VAL-01 and VAL-04 may present agreed synthetic scenarios to actual stakeholder representatives to assess their decisions, with the simulation limitation recorded. Participant feedback and technical acceptance remain Pending. Environment/contract approvals: TBD; owner Jack Marino; decision date TBD.
\subsection{Main Test Procedure Readiness}
\label{Section:MainTestProcedures}
The 12 procedures below are prepared for future execution. No hardware is
available yet; no test has been run. Every test result and verification result
is \StatusPending; evidence remains Not collected. An unresolved setup value
blocks execution or acceptance, not preparation of the procedure.
All TBD values, including hardware, input counts, durations, instrumentation,
software choices and execution dates, are owned by Jack Marino; decision date:
TBD. Each procedure names its applicable parameter decisions. The 10 ms and
100 ms values are owner-selected testing targets, with final approval pending
results; no other numerical bound is invented. OMS remains Proposed.

\paragraph{Common preparation and run record.}
Before a run, approve the applicable parameter values and setup, identify the
requirement/check and product/source revision, hardware/firmware/driver/host
versions, instrument calibration and uncertainty, configuration, operator,
execution date, selected inputs, counts and duration. All are TBD until supplied.
Prepare the expected observations before collecting data. Log each procedure
step, raw observations, calculations, deviations and evidence location.
The mapped verification entry is the authoritative acceptance bound.
Only an executed, reviewed run may change a result; incomplete setup or
ambiguous evidence leaves it Pending with a recorded blocker. Record the
actual outcome, including any failure, rather than copying an expected result.
Restore the baseline configuration after each test.

\paragraph{Coverage and timing boundaries.}\label{Section:TestReviewClarifications}
For acquisition and latency tests, the independent conversion-ready trace defines the event population; software acceptance logs cannot remove missed events from the denominator. The workload/schedule and treatment of assertions outside it remain TBD under \hyperref[Parameter:PAR-01]{PAR-01}, \hyperref[Parameter:PAR-03]{PAR-03}, and \hyperref[Parameter:PAR-04]{PAR-04}. The unchanged lecture timing statement is not interpreted as a demonstrated guarantee under unlimited overload. Overload behavior is separately checked by \TestCaseRef{TC-026}. No choice of timestamp error, sample rate, buffer size or uncertainty guard band is made by this review. A missing event, measurement or unresolved acceptance dependency cannot produce a Pass. Expected results below summarize the mapped verification entry; that entry remains the acceptance authority.

\paragraph{Software-only rehearsal.}
Saved files or controlled test inputs may help rehearse recording, replay,
record decoding, access checks or gateway mappings after those interfaces
exist. Label such evidence with its actual scope. A rehearsal does not establish
physical ADC capture, hardware timing, disconnect recovery or platform/hardware
support. No rehearsal has been executed in this update.
\begin{landscape}
\small
\begin{tabularx}{\linewidth}{L{0.8cm}L{2.0cm}Y L{2.3cm}L{2.3cm}L{2.1cm}}
\caption{Main Baseline Test Procedure Index\label{Table:WeekTwoTests}}\\
\toprule
\textbf{No.} & \textbf{Procedure} & \textbf{Requirement} & \textbf{Verification} & \textbf{Planned evidence} & \textbf{Result}\\
\midrule
\endfirsthead
\textbf{No.} & \textbf{Procedure} & \textbf{Requirement} & \textbf{Verification} & \textbf{Planned evidence} & \textbf{Result}\\
\midrule
\endhead
\midrule
\multicolumn{6}{r}{\small\itshape Continued on next page}\\
\endfoot
\bottomrule
\endlastfoot
1 & \TestCaseRef{TC-010} & \RequirementRef{REQ-SMP-01} & \VerificationRef{VER-010} & \ArtifactRef{Evidence}{EVD-010} & \StatusPending\\
2 & \TestCaseRef{TC-016} & \RequirementRef{REQ-IF-01} & \VerificationRef{VER-016} & \ArtifactRef{Evidence}{EVD-016} & \StatusPending\\
3 & \TestCaseRef{TC-019} & \RequirementRef{REQ-TIM-01} & \VerificationRef{VER-019} & \ArtifactRef{Evidence}{EVD-019} & \StatusPending\\
4 & \TestCaseRef{TC-027} & \RequirementRef{REQ-SMP-02} & \VerificationRef{VER-027} & \ArtifactRef{Evidence}{EVD-027} & \StatusPending\\
5 & \TestCaseRef{TC-026} & \RequirementRef{REQ-F-009} & \VerificationRef{VER-026} & \ArtifactRef{Evidence}{EVD-026} & \StatusPending\\
6 & \TestCaseRef{TC-050} & \RequirementRef{REQ-F-018} & \VerificationRef{VER-050} & \ArtifactRef{Evidence}{EVD-050} & \StatusPending\\
7 & \TestCaseRef{TC-051} & \RequirementRef{REQ-F-019} & \VerificationRef{VER-051} & \ArtifactRef{Evidence}{EVD-051} & \StatusPending\\
8 & \TestCaseRef{TC-029} & \RequirementRef{REQ-F-010} & \VerificationRef{VER-029} & \ArtifactRef{Evidence}{EVD-029} & \StatusPending\\
9 & \TestCaseRef{TC-052} & \RequirementRef{REQ-Q-016} & \VerificationRef{VER-052} & \ArtifactRef{Evidence}{EVD-052} & \StatusPending\\
10 & \TestCaseRef{TC-053} & \RequirementRef{REQ-S-002} & \VerificationRef{VER-053} & \ArtifactRef{Evidence}{EVD-053} & \StatusPending\\
11 & \TestCaseRef{TC-032} & \RequirementRef{REQ-OMS-01} & \VerificationRef{VER-032} & \ArtifactRef{Evidence}{EVD-032} & \StatusPending\\
12 & \TestCaseRef{TC-054} & \RequirementRef{REQ-I-006} & \VerificationRef{VER-054} & \ArtifactRef{Evidence}{EVD-054} & \StatusPending\\
\end{tabularx}
\normalsize
\end{landscape}
\subsection{1. Completed ADC reading}
\small
\begin{tabularx}{\linewidth}{L{3.2cm}Y}
\caption{Completed ADC reading\label{Table:Procedure010}}\\
\toprule
\textbf{Field / Step} & \textbf{Planned procedure}\\
\midrule
\endfirsthead
\textbf{Field / Step} & \textbf{Planned procedure}\\
\midrule
\endhead
\midrule
\multicolumn{2}{r}{\small\itshape Continued on next page}\\
\endfoot
\bottomrule
\endlastfoot
Identifier & \TestCaseID{TC-010}\\
Coverage & \RequirementRef{REQ-SMP-01}; \UseCaseRef{UC-01}; \VerificationRef{VER-010}\\
Result & \StatusPending. Not executed; actual observations, execution date and evidence location: TBD.\\
Owner / Decision date & Jack Marino / TBD\\
Entry conditions & Resolve applicable \hyperref[Parameter:PAR-01]{PAR-01}, \hyperref[Parameter:PAR-02]{PAR-02} decisions and the \hyperref[Section:MainTestProcedures]{common preparation}.\\
Setup / hardware TBDs & Hardware boards/revisions, voltage/reference range, ADC clock, conversion-ready observation, control wiring, FT245R mode and independent bus-capture instrument: TBD.\\
Inputs / test cases & Selected potentiometer positions within the approved voltage range; independently observed completed-conversion codes and ready events. Positions, repetition count and run duration: TBD.\\
Step 1 & Record the approved hardware setup and start independent capture of conversion-ready events, the full ADC output code and host publications.\\
Step 2 & Set the potentiometer to each selected position and allow the approved settling time (TBD). Initiate conversions using the selected control method (TBD).\\
Step 3 & Enumerate every conversion-ready assertion in the approved acquisition schedule from the independent trace, not from the host acceptance log. Read its completed 8-bit code and locate its corresponding host publication.\\
Step 4 & Compare the published code with the independently captured code and inspect the published validity state. Record every missing, duplicate, mismatched or unqualified publication.\\
Step 5 & Stop acquisition and retain the event-to-record comparison. Do not infer an exact ADC code from knob position alone.\\
Expected result & Every conversion-ready assertion in the approved acquisition schedule has a matching publication with the independently observed 8-bit code and its validity state. Missing or ambiguous matches prevent a pass disposition. Authoritative bound: \VerificationRef{VER-010}. This is an expectation, not an observed result.\\
Planned evidence & \ArtifactRef{Evidence}{EVD-010} (Not collected). Ready-event/bus trace, published records and per-event code/validity reconciliation.\\
Record / restore & Record per-step observations and any deviation against the expected result, retain raw data and calculations, restore the baseline, and leave the result Pending until an executed review supplies a disposition.\\
Navigation & \hyperref[Table:WeekTwoTests]{Main procedure index}; \hyperref[Table:TraceabilityMatrix]{main trace matrix}.\\
\end{tabularx}
\normalsize
\subsection{2. Local sample record fields}
\small
\begin{tabularx}{\linewidth}{L{3.2cm}Y}
\caption{Local sample record fields\label{Table:Procedure016}}\\
\toprule
\textbf{Field / Step} & \textbf{Planned procedure}\\
\midrule
\endfirsthead
\textbf{Field / Step} & \textbf{Planned procedure}\\
\midrule
\endhead
\midrule
\multicolumn{2}{r}{\small\itshape Continued on next page}\\
\endfoot
\bottomrule
\endlastfoot
Identifier & \TestCaseID{TC-016}\\
Coverage & \RequirementRef{REQ-IF-01}; \UseCaseRef{UC-01}; \VerificationRef{VER-016}\\
Result & \StatusPending. Not executed; actual observations, execution date and evidence location: TBD.\\
Owner / Decision date & Jack Marino / TBD\\
Entry conditions & Resolve applicable \hyperref[Parameter:PAR-02]{PAR-02}, \hyperref[Parameter:PAR-03]{PAR-03} decisions and the \hyperref[Section:MainTestProcedures]{common preparation}.\\
Setup / hardware TBDs & Approved local schema, validity states, timestamp event/units/epoch/clock, sequence width, acquisition-epoch rules and independent decoder: TBD. Hardware capture configuration and input rate: TBD.\\
Inputs / test cases & A controlled sequence of qualified readings with distinct code changes; run length and sequence-boundary cases: TBD.\\
Step 1 & Record the schema and clock/sequence conventions; start independent capture of the input sequence and emitted local records.\\
Step 2 & Acquire the selected sequence and decode every emitted record with the agreed independent decoder.\\
Step 3 & Check that each qualified record contains code, validity, timestamp and sequence number with the approved types and encoding; compare the code and validity with the input evidence.\\
Step 4 & Check timestamps against the declared event and clock convention; check increasing sequence order within the defined epoch. If a restart or boundary is included, evaluate it under the approved epoch policy.\\
Step 5 & List missing fields, decoding errors, sequence duplicates/regressions and unexplained gaps. Retain the decoded records and report.\\
Expected result & Every qualified record contains all four required fields with the approved types and encoding; the code is an unsigned 8-bit value and sequence numbers strictly increase within the declared acquisition epoch. Restart and wrap follow the approved epoch policy. Authoritative bound: \VerificationRef{VER-016}. This is an expectation, not an observed result.\\
Planned evidence & \ArtifactRef{Evidence}{EVD-016} (Not collected). Schema revision, raw and decoded records, per-field checks and epoch/sequence-order report. No supporting 1 ms error bound is assumed.\\
Record / restore & Record per-step observations and any deviation against the expected result, retain raw data and calculations, restore the baseline, and leave the result Pending until an executed review supplies a disposition.\\
Navigation & \hyperref[Table:WeekTwoTests]{Main procedure index}; \hyperref[Table:TraceabilityMatrix]{main trace matrix}.\\
\end{tabularx}
\normalsize
\subsection{3. Publication latency: 10 ms testing target}
\small
\begin{tabularx}{\linewidth}{L{3.2cm}Y}
\caption{Publication latency: 10 ms testing target\label{Table:Procedure019}}\\
\toprule
\textbf{Field / Step} & \textbf{Planned procedure}\\
\midrule
\endfirsthead
\textbf{Field / Step} & \textbf{Planned procedure}\\
\midrule
\endhead
\midrule
\multicolumn{2}{r}{\small\itshape Continued on next page}\\
\endfoot
\bottomrule
\endlastfoot
Identifier & \TestCaseID{TC-019}\\
Coverage & \RequirementRef{REQ-TIM-01}; \UseCaseRef{UC-01}; \VerificationRef{VER-019}\\
Result & \StatusPending. Not executed; actual observations, execution date and evidence location: TBD.\\
Owner / Decision date & Jack Marino / TBD\\
Entry conditions & Resolve applicable \hyperref[Parameter:PAR-01]{PAR-01}, \hyperref[Parameter:PAR-02]{PAR-02}, \hyperref[Parameter:PAR-03]{PAR-03}, \hyperref[Parameter:PAR-04]{PAR-04} decisions and the \hyperref[Section:MainTestProcedures]{common preparation}.\\
Setup / hardware TBDs & Physical ready-event capture, observable publication boundary, common or correlated clocks, reference calibration, uncertainty treatment, load envelope and acquisition hardware: TBD.\\
Inputs / test cases & Completed conversions at the approved normal and boundary workloads. Rates, burst patterns, duration and number of events: TBD.\\
Step 1 & Record the ready-event definition and publication boundary. Establish a common or correlated measurement time basis and its uncertainty before acquisition.\\
Step 2 & Apply each approved workload while capturing ready-event times and corresponding local-publication times.\\
Step 3 & Enumerate every ready assertion from the independent trace within each approved workload case and pair it with exactly one corresponding publication. Record missing, duplicate or ambiguous matches; none may be omitted or counted as a successful timing case.\\
Step 4 & For each pair calculate publication time minus ready-event time, applying the approved clock correlation and uncertainty treatment.\\
Step 5 & Compare every result against the 10 ms testing target and report the maximum, any target exceedance and all unmatched events. Averages alone do not establish compliance.\\
Expected result & Every ready assertion in the approved workload case has exactly one corresponding publication, and every latency is at most 10 ms under the agreed uncertainty treatment. Missing, duplicate or ambiguous matches prevent a pass disposition. Authoritative bound: \VerificationRef{VER-019}. This is an expectation, not an observed result.\\
Planned evidence & \ArtifactRef{Evidence}{EVD-019} (Not collected). Paired event/publication timestamps, workload configuration, clock-correlation/calibration evidence, per-event latency and uncertainty calculations, maximum and exception list.\\
Record / restore & Record per-step observations and any deviation against the expected result, retain raw data and calculations, restore the baseline, and leave the result Pending until an executed review supplies a disposition.\\
Navigation & \hyperref[Table:WeekTwoTests]{Main procedure index}; \hyperref[Table:TraceabilityMatrix]{main trace matrix}.\\
\end{tabularx}
\normalsize
\subsection{4. Stale data: 100 ms testing target}
\small
\begin{tabularx}{\linewidth}{L{3.2cm}Y}
\caption{Stale data: 100 ms testing target\label{Table:Procedure027}}\\
\toprule
\textbf{Field / Step} & \textbf{Planned procedure}\\
\midrule
\endfirsthead
\textbf{Field / Step} & \textbf{Planned procedure}\\
\midrule
\endhead
\midrule
\multicolumn{2}{r}{\small\itshape Continued on next page}\\
\endfoot
\bottomrule
\endlastfoot
Identifier & \TestCaseID{TC-027}\\
Coverage & \RequirementRef{REQ-SMP-02}; \UseCaseRef{UC-02}; \VerificationRef{VER-027}\\
Result & \StatusPending. Not executed; actual observations, execution date and evidence location: TBD.\\
Owner / Decision date & Jack Marino / TBD\\
Entry conditions & Resolve applicable \hyperref[Parameter:PAR-01]{PAR-01}, \hyperref[Parameter:PAR-03]{PAR-03}, \hyperref[Parameter:PAR-06]{PAR-06} decisions and the \hyperref[Section:MainTestProcedures]{common preparation}.\\
Setup / hardware TBDs & Controlled method to stop new qualified samples, observable stale-status boundary, clock/reference and uncertainty treatment: TBD. Hardware wiring, sample interval, observation resolution and repetition count: TBD.\\
Inputs / test cases & A valid publication followed by a controlled pause in qualified input; a restoration of input after stale status.\\
Step 1 & Start acquisition and capture all valid publications and status transitions using the approved reference; retain their sample identities.\\
Step 2 & Stop new qualified input using the approved method while keeping the status observer active. Capture any in-flight or buffered samples that still publish; identify the actual last valid publication after these publications, using its identity and time as the stale-time reference.\\
Step 3 & Capture the stale-state transition attributable to the absence of new qualified input after that actual last valid publication. Calculate their elapsed time with the approved uncertainty treatment; do not start the stale timer at the pause command or USB disconnect time.\\
Step 4 & Compare the stale-state transition against the 100 ms testing target; inspect any retained reading to ensure it is not presented as newly valid.\\
Step 5 & Restore qualified input and record the subsequent state. Repeat with the approved pause positions/counts (TBD); retain late, missing and ambiguous transitions.\\
Expected result & The current acquisition output becomes stale within 100 ms of the previous valid publication; no new sample is fabricated. Authoritative bound: \VerificationRef{VER-027}. This is an expectation, not an observed result.\\
Planned evidence & \ArtifactRef{Evidence}{EVD-027} (Not collected). Last-valid-publication time, controlled pause trace, stale-status transition time, elapsed-time/uncertainty calculation and retained-value status observations.\\
Record / restore & Record per-step observations and any deviation against the expected result, retain raw data and calculations, restore the baseline, and leave the result Pending until an executed review supplies a disposition.\\
Navigation & \hyperref[Table:WeekTwoTests]{Main procedure index}; \hyperref[Table:TraceabilityMatrix]{main trace matrix}.\\
\end{tabularx}
\normalsize
\subsection{5. Visible overload and lost readings}
\small
\begin{tabularx}{\linewidth}{L{3.2cm}Y}
\caption{Visible overload and lost readings\label{Table:Procedure026}}\\
\toprule
\textbf{Field / Step} & \textbf{Planned procedure}\\
\midrule
\endfirsthead
\textbf{Field / Step} & \textbf{Planned procedure}\\
\midrule
\endhead
\midrule
\multicolumn{2}{r}{\small\itshape Continued on next page}\\
\endfoot
\bottomrule
\endlastfoot
Identifier & \TestCaseID{TC-026}\\
Coverage & \RequirementRef{REQ-F-009}; \UseCaseRef{UC-02}; \VerificationRef{VER-026}\\
Result & \StatusPending. Not executed; actual observations, execution date and evidence location: TBD.\\
Owner / Decision date & Jack Marino / TBD\\
Entry conditions & Resolve applicable \hyperref[Parameter:PAR-02]{PAR-02}, \hyperref[Parameter:PAR-04]{PAR-04}, \hyperref[Parameter:PAR-06]{PAR-06} decisions and the \hyperref[Section:MainTestProcedures]{common preparation}.\\
Setup / hardware TBDs & Input/service envelope, finite buffers, overflow/backpressure/drop policy, loss/discontinuity interface and queue instrumentation: TBD. Hardware source rate and workload controls: TBD.\\
Inputs / test cases & Identifiable sample streams at an approved normal load and loads that force the selected overload condition; rates, stall lengths, duration and expected accounting rules: TBD.\\
Step 1 & Record the approved overload policy and initial queue state. Run the normal-load case and retain source, accepted, delivered and queued-record counts.\\
Step 2 & Increase input load or stall the consumer using the approved mechanism until the configured overload condition is reached.\\
Step 3 & Observe whether the system follows its selected backpressure or drop policy; capture queue occupancy and any loss/discontinuity indication seen by the receiving software.\\
Step 4 & Resume the consumer and reconcile input, delivered, retained and declared lost or uncertain records under the approved accounting rules.\\
Step 5 & Report any silent loss, unexplained gap or presentation of interrupted acquisition as continuous. Do not assume a drop policy or zero loss before the policy is selected.\\
Expected result & The selected overflow policy is followed; loss or uncertainty is visible and is never represented as uninterrupted valid acquisition. Authoritative bound: \VerificationRef{VER-026}. This is an expectation, not an observed result.\\
Planned evidence & \ArtifactRef{Evidence}{EVD-026} (Not collected). Source/consumer counts, queue trace, selected policy, overload trigger, backpressure/drop observations and consumer-visible loss/discontinuity accounting.\\
Record / restore & Record per-step observations and any deviation against the expected result, retain raw data and calculations, restore the baseline, and leave the result Pending until an executed review supplies a disposition.\\
Navigation & \hyperref[Table:WeekTwoTests]{Main procedure index}; \hyperref[Table:TraceabilityMatrix]{main trace matrix}.\\
\end{tabularx}
\normalsize
\subsection{6. Save readings}
\small
\begin{tabularx}{\linewidth}{L{3.2cm}Y}
\caption{Save readings\label{Table:Procedure050}}\\
\toprule
\textbf{Field / Step} & \textbf{Planned procedure}\\
\midrule
\endfirsthead
\textbf{Field / Step} & \textbf{Planned procedure}\\
\midrule
\endhead
\midrule
\multicolumn{2}{r}{\small\itshape Continued on next page}\\
\endfoot
\bottomrule
\endlastfoot
Identifier & \TestCaseID{TC-050}\\
Coverage & \RequirementRef{REQ-F-018}; \UseCaseRef{UC-06}; \VerificationRef{VER-050}\\
Result & \StatusPending. Not executed; actual observations, execution date and evidence location: TBD.\\
Owner / Decision date & Jack Marino / TBD\\
Entry conditions & Resolve applicable \hyperref[Parameter:PAR-02]{PAR-02}, \hyperref[Parameter:PAR-13]{PAR-13} decisions and the \hyperref[Section:MainTestProcedures]{common preparation}.\\
Setup / hardware TBDs & JSON Lines and basic recording controls approved. Detailed selection semantics/schema, metadata encoding, storage/retention limits, independent reader and storage-failure behavior: TBD. Real acquisition hardware and rate: TBD; a software-only rehearsal must be labeled separately.\\
Inputs / test cases & A finite identifiable stream containing valid and invalid records, gaps and acquisition-epoch metadata. Dataset length and safe isolated storage-failure mechanism: TBD.\\
Step 1 & Retain the reference stream and its metadata; select the records to save according to the approved recording controls.\\
Step 2 & Start recording, submit the stream, stop recording and close the dataset using the documented workflow.\\
Step 3 & Reopen the saved dataset with the independent reader. Compare each accepted record, original order, code, validity, timestamp, sequence and epoch context with the reference.\\
Step 4 & In an isolated test destination, trigger the approved storage-failure case; record which samples were saved, lost or uncertain and the visible failure indication.\\
Step 5 & Preserve both normal and failure datasets and the reconciliation report; do not alter the reference data.\\
Expected result & Every record accepted for recording matches its source fields and order; storage failures and any recording loss are visible under the approved recording policy. Authoritative bound: \VerificationRef{VER-050}. This is an expectation, not an observed result.\\
Planned evidence & \ArtifactRef{Evidence}{EVD-050} (Not collected). Reference stream, selected-record list, file/schema metadata, normal and failure datasets, field/order reconciliation and storage-failure indication.\\
Record / restore & Record per-step observations and any deviation against the expected result, retain raw data and calculations, restore the baseline, and leave the result Pending until an executed review supplies a disposition.\\
Navigation & \hyperref[Table:WeekTwoTests]{Main procedure index}; \hyperref[Table:TraceabilityMatrix]{main trace matrix}.\\
\end{tabularx}
\normalsize
\subsection{7. Replay saved readings}
\small
\begin{tabularx}{\linewidth}{L{3.2cm}Y}
\caption{Replay saved readings\label{Table:Procedure051}}\\
\toprule
\textbf{Field / Step} & \textbf{Planned procedure}\\
\midrule
\endfirsthead
\textbf{Field / Step} & \textbf{Planned procedure}\\
\midrule
\endhead
\midrule
\multicolumn{2}{r}{\small\itshape Continued on next page}\\
\endfoot
\bottomrule
\endlastfoot
Identifier & \TestCaseID{TC-051}\\
Coverage & \RequirementRef{REQ-F-019}; \UseCaseRef{UC-06}; \VerificationRef{VER-051}\\
Result & \StatusPending. Not executed; actual observations, execution date and evidence location: TBD.\\
Owner / Decision date & Jack Marino / TBD\\
Entry conditions & Resolve applicable \hyperref[Parameter:PAR-02]{PAR-02}, \hyperref[Parameter:PAR-13]{PAR-13} decisions and the \hyperref[Section:MainTestProcedures]{common preparation}.\\
Setup / hardware TBDs & JSON Lines, basic playback controls and explicit Replay mode approved. Detailed dataset schema, reader compatibility, pacing, replay-output boundary and indicator representation: TBD. Hardware is not required for a file-only rehearsal, but that rehearsal is not evidence of physical acquisition.\\
Inputs / test cases & A known saved dataset with valid/invalid records, explicit gaps and original timestamps, sequences and epoch context; selection and playback settings: TBD.\\
Step 1 & Independently decode and retain the expected selected dataset contents before starting playback.\\
Step 2 & Open the dataset and start replay with the approved controls and pacing; capture all consumer-visible output and its replay indication.\\
Step 3 & Compare the replayed records field by field with the saved originals and verify their selected order. Record omissions, duplicates, metadata changes and invalid-state changes.\\
Step 4 & Check that each replay output is identifiable as recorded data, including around the start/end and any approved pause/resume operation; original timestamps remain original.\\
Step 5 & Compare playback timing with the selected pacing contract without inventing a speed requirement; retain the input/output comparison and mode observations.\\
Expected result & All selected replay records match saved metadata and order, and every replay output is distinguishable from live output; pacing follows the approved replay contract. Authoritative bound: \VerificationRef{VER-051}. This is an expectation, not an observed result.\\
Planned evidence & \ArtifactRef{Evidence}{EVD-051} (Not collected). Saved input, independent decoded expectation, replay output, playback-control/pacing settings, field/order comparison and replay-mode indication.\\
Record / restore & Record per-step observations and any deviation against the expected result, retain raw data and calculations, restore the baseline, and leave the result Pending until an executed review supplies a disposition.\\
Navigation & \hyperref[Table:WeekTwoTests]{Main procedure index}; \hyperref[Table:TraceabilityMatrix]{main trace matrix}.\\
\end{tabularx}
\normalsize
\subsection{8. Qualified recovery after disconnect or reset}
\small
\begin{tabularx}{\linewidth}{L{3.2cm}Y}
\caption{Qualified recovery after disconnect or reset\label{Table:Procedure029}}\\
\toprule
\textbf{Field / Step} & \textbf{Planned procedure}\\
\midrule
\endfirsthead
\textbf{Field / Step} & \textbf{Planned procedure}\\
\midrule
\endhead
\midrule
\multicolumn{2}{r}{\small\itshape Continued on next page}\\
\endfoot
\bottomrule
\endlastfoot
Identifier & \TestCaseID{TC-029}\\
Coverage & \RequirementRef{REQ-F-010}; \UseCaseRef{UC-02}; \VerificationRef{VER-029}\\
Result & \StatusPending. Not executed; actual observations, execution date and evidence location: TBD.\\
Owner / Decision date & Jack Marino / TBD\\
Entry conditions & Resolve applicable \hyperref[Parameter:PAR-01]{PAR-01}, \hyperref[Parameter:PAR-02]{PAR-02}, \hyperref[Parameter:PAR-03]{PAR-03}, \hyperref[Parameter:PAR-06]{PAR-06} decisions and the \hyperref[Section:MainTestProcedures]{common preparation}.\\
Setup / hardware TBDs & Device identity checks, configuration baseline, timing qualification, restart epoch/queue policy and approved disconnect/reset injection: TBD. Hardware/driver behavior and recovery observation points: TBD.\\
Inputs / test cases & An active valid stream followed by disconnect and reset cases; withheld identity/configuration/timing/new-sample prerequisites and controlled recovery sequence: TBD.\\
Step 1 & Acquire valid readings and retain the device identity, active configuration, timing state and acquisition epoch.\\
Step 2 & Apply one approved disconnect or reset. Capture status, queued-data handling and every subsequent publication.\\
Step 3 & Reconnect or restore the path while deliberately withholding each recovery prerequisite in separate approved cases; verify that valid output remains withheld.\\
Step 4 & Restore correct identity, configuration and timing qualification, then permit a newly completed sample. Record the ordering of prerequisite completion and the first valid publication.\\
Step 5 & Check that obsolete queued samples are not presented as new. Repeat the selected reset/disconnect cases and retain failed prerequisites; no recovery deadline is invented.\\
Expected result & Valid output resumes only after all approved recovery predicates and a new qualified sample; obsolete queued samples are not presented as new. Authoritative bound: \VerificationRef{VER-029}. This is an expectation, not an observed result.\\
Planned evidence & \ArtifactRef{Evidence}{EVD-029} (Not collected). Fault-injection timeline, device/configuration/timing checks, epoch and queue disposition, withheld-prerequisite cases and first-new-sample/valid-output ordering.\\
Record / restore & Record per-step observations and any deviation against the expected result, retain raw data and calculations, restore the baseline, and leave the result Pending until an executed review supplies a disposition.\\
Navigation & \hyperref[Table:WeekTwoTests]{Main procedure index}; \hyperref[Table:TraceabilityMatrix]{main trace matrix}.\\
\end{tabularx}
\normalsize
\subsection{9. Supported macOS, Windows and Linux workflow}
\small
\begin{tabularx}{\linewidth}{L{3.2cm}Y}
\caption{Supported macOS, Windows and Linux workflow\label{Table:Procedure052}}\\
\toprule
\textbf{Field / Step} & \textbf{Planned procedure}\\
\midrule
\endfirsthead
\textbf{Field / Step} & \textbf{Planned procedure}\\
\midrule
\endhead
\midrule
\multicolumn{2}{r}{\small\itshape Continued on next page}\\
\endfoot
\bottomrule
\endlastfoot
Identifier & \TestCaseID{TC-052}\\
Coverage & \RequirementRef{REQ-Q-016}; \UseCaseRef{UC-01}; \VerificationRef{VER-052}\\
Result & \StatusPending. Not executed; actual observations, execution date and evidence location: TBD.\\
Owner / Decision date & Jack Marino / TBD\\
Entry conditions & Resolve applicable \hyperref[Parameter:PAR-01]{PAR-01}, \hyperref[Parameter:PAR-12]{PAR-12} decisions and the \hyperref[Section:MainTestProcedures]{common preparation}.\\
Setup / hardware TBDs & Supported OS/CPU/driver matrix, host resources/privileges, deployment package and instructions, acquisition settings and hardware: TBD. No OS version is assumed supported before approval.\\
Inputs / test cases & One agreed installation, configuration, acquisition and status-inspection workflow on every approved matrix entry; input sequence, duration and dependency-failure cases: TBD.\\
Step 1 & Freeze and record the common application source revision and approved platform/driver matrix.\\
Step 2 & On each matrix entry, follow the documented installation and configuration instructions; record versions, resources, permissions and any deviations.\\
Step 3 & Connect the selected hardware and run the agreed acquisition/status workflow with the controlled input; retain decoded readings and visible status.\\
Step 4 & Compare the source revision across platforms and identify any platform-specific application source edits. Record allowed platform-specific packaging/configuration separately.\\
Step 5 & Check that missing prerequisites can be identified from deployment instructions. Report success or obstruction for every matrix entry; untested entries remain Pending.\\
Expected result & Every approved matrix entry completes the agreed workflow without platform-specific application source changes; unmet prerequisites are identified in deployment instructions. Authoritative bound: \VerificationRef{VER-052}. This is an expectation, not an observed result.\\
Planned evidence & \ArtifactRef{Evidence}{EVD-052} (Not collected). Approved matrix, common source revision, per-platform installation/configuration/acquisition/status records, driver dependencies, deviations and per-entry disposition.\\
Record / restore & Record per-step observations and any deviation against the expected result, retain raw data and calculations, restore the baseline, and leave the result Pending until an executed review supplies a disposition.\\
Navigation & \hyperref[Table:WeekTwoTests]{Main procedure index}; \hyperref[Table:TraceabilityMatrix]{main trace matrix}.\\
\end{tabularx}
\normalsize
\subsection{10. Authorized acquisition-setting changes}
\small
\begin{tabularx}{\linewidth}{L{3.2cm}Y}
\caption{Authorized acquisition-setting changes\label{Table:Procedure053}}\\
\toprule
\textbf{Field / Step} & \textbf{Planned procedure}\\
\midrule
\endfirsthead
\textbf{Field / Step} & \textbf{Planned procedure}\\
\midrule
\endhead
\midrule
\multicolumn{2}{r}{\small\itshape Continued on next page}\\
\endfoot
\bottomrule
\endlastfoot
Identifier & \TestCaseID{TC-053}\\
Coverage & \RequirementRef{REQ-S-002}; \UseCaseRef{UC-01}; \VerificationRef{VER-053}\\
Result & \StatusPending. Not executed; actual observations, execution date and evidence location: TBD.\\
Owner / Decision date & Jack Marino / TBD\\
Entry conditions & Resolve applicable \hyperref[Parameter:PAR-14]{PAR-14} decisions and the \hyperref[Section:MainTestProcedures]{common preparation}.\\
Setup / hardware TBDs & Protected setting/operation list, identity/access policy, enforcement boundary, interfaces and denial reporting: TBD. Hardware-dependent setting limits and acquisition state: TBD.\\
Inputs / test cases & For every protected operation: approved changes attempted by authorized, unauthorized and missing/unauthenticated identities. Test identities and values: TBD.\\
Step 1 & Record the policy revision, protected-operation list and baseline effective settings. Build the identity-by-operation test matrix before executing requests.\\
Step 2 & For each authorized case, request a permitted change through the actual enforcement interface and independently inspect the effective setting.\\
Step 3 & Restore the baseline; attempt the corresponding change with unauthorized and missing identities through each applicable protected interface.\\
Step 4 & Compare effective settings before and after each denied request and capture the required denial indication. Do not infer denial from a screen message alone.\\
Step 5 & Reconcile all cases against the approved permission matrix, retain unexpected grants/denials and restore the baseline configuration.\\
Expected result & Authorized changes follow the approved policy; every unauthorized or unauthenticated protected change is denied and leaves effective settings unchanged, with the required denial indication. Authoritative bound: \VerificationRef{VER-053}. This is an expectation, not an observed result.\\
Planned evidence & \ArtifactRef{Evidence}{EVD-053} (Not collected). Policy and identity/operation matrix, request/response observations, independently inspected before/after effective settings, denial indications and restoration record.\\
Record / restore & Record per-step observations and any deviation against the expected result, retain raw data and calculations, restore the baseline, and leave the result Pending until an executed review supplies a disposition.\\
Navigation & \hyperref[Table:WeekTwoTests]{Main procedure index}; \hyperref[Table:TraceabilityMatrix]{main trace matrix}.\\
\end{tabularx}
\normalsize
\subsection{11. Proposed OMS sensor-status exchange}
\small
\begin{tabularx}{\linewidth}{L{3.2cm}Y}
\caption{Proposed OMS sensor-status exchange\label{Table:Procedure032}}\\
\toprule
\textbf{Field / Step} & \textbf{Planned procedure}\\
\midrule
\endfirsthead
\textbf{Field / Step} & \textbf{Planned procedure}\\
\midrule
\endhead
\midrule
\multicolumn{2}{r}{\small\itshape Continued on next page}\\
\endfoot
\bottomrule
\endlastfoot
Identifier & \TestCaseID{TC-032}\\
Coverage & \RequirementRef{REQ-OMS-01}; \UseCaseRef{UC-03}; \VerificationRef{VER-032}\\
Result & \StatusPending. Not executed; actual observations, execution date and evidence location: TBD.\\
Owner / Decision date & Jack Marino / TBD\\
Entry conditions & Resolve applicable \hyperref[Parameter:PAR-07]{PAR-07} decisions and the \hyperref[Section:MainTestProcedures]{common preparation}.\\
Setup / hardware TBDs & OMS applicability, standard version, CAL profile, UCI message, status mapping, responsible gateway, service/failure contract and independent decoder: TBD. Interface test harness, hardware source and timing allocations: TBD. Requirement remains Proposed; execution is gated.\\
Inputs / test cases & Approved valid, stale and faulted local states, with correlated source records and selected mapping/failure cases. Exact messages, repetitions and test endpoint: TBD.\\
Step 1 & Check the architecture decision and approval record first. If any required OMS decision is unresolved, record the blocker and leave this test Pending without claiming OMS conformance.\\
Step 2 & After approval, configure the selected gateway and independent observer at the defined CAL boundary; retain contract/profile/message revisions.\\
Step 3 & Supply each approved valid, stale and faulted local state and capture the corresponding mission-facing product.\\
Step 4 & Decode each product independently and compare fields, state mapping and boundary behavior against the approved UCI/service contract; correlate products with their source states.\\
Step 5 & Exercise the approved mapping/delivery failure cases and retain explicit fault observations. Report mismatches without substituting the local sample format for OMS evidence.\\
Expected result & Each state produces the approved derived status product at the CAL boundary using the approved mapping; local sample records are not treated as the OMS contract. Authoritative bound: \VerificationRef{VER-032}. This is an expectation, not an observed result.\\
Planned evidence & \ArtifactRef{Evidence}{EVD-032} (Not collected). OMS applicability/approval decision, version/profile/message/mapping baseline, gateway identity, injected states, decoded CAL captures, correlation and contract-fault observations.\\
Record / restore & Record per-step observations and any deviation against the expected result, retain raw data and calculations, restore the baseline, and leave the result Pending until an executed review supplies a disposition.\\
Navigation & \hyperref[Table:WeekTwoTests]{Main procedure index}; \hyperref[Table:TraceabilityMatrix]{main trace matrix}.\\
\end{tabularx}
\normalsize
\subsection{12. Mission timestamp and timing quality}
\small
\begin{tabularx}{\linewidth}{L{3.2cm}Y}
\caption{Mission timestamp and timing quality\label{Table:Procedure054}}\\
\toprule
\textbf{Field / Step} & \textbf{Planned procedure}\\
\midrule
\endfirsthead
\textbf{Field / Step} & \textbf{Planned procedure}\\
\midrule
\endhead
\midrule
\multicolumn{2}{r}{\small\itshape Continued on next page}\\
\endfoot
\bottomrule
\endlastfoot
Identifier & \TestCaseID{TC-054}\\
Coverage & \RequirementRef{REQ-I-006}; \UseCaseRef{UC-03}; \VerificationRef{VER-054}\\
Result & \StatusPending. Not executed; actual observations, execution date and evidence location: TBD.\\
Owner / Decision date & Jack Marino / TBD\\
Entry conditions & Resolve applicable \hyperref[Parameter:PAR-02]{PAR-02}, \hyperref[Parameter:PAR-03]{PAR-03}, \hyperref[Parameter:PAR-07]{PAR-07} decisions and the \hyperref[Section:MainTestProcedures]{common preparation}.\\
Setup / hardware TBDs & Approved mission contract/gateway, timestamp event/units/epoch, timing-quality states, unavailable-time representation and clock correlation: TBD. Hardware capture/reference and source timestamps: TBD; OMS-specific execution depends on its decision.\\
Inputs / test cases & Sample-related status for qualified, unqualified and unavailable timing, plus repeated or delayed status for identifiable source samples. Exact cases and delays: TBD.\\
Step 1 & Record the approved source-to-mission timestamp/quality mapping and prepare independently identifiable samples covering every agreed timing-quality state.\\
Step 2 & Submit the corresponding local sample/status cases to the gateway and capture mission-facing status alongside original source records.\\
Step 3 & Correlate each sample-related product with its source. Decode its timestamp event, units and epoch; verify semantic preservation under the approved representation, not replacement by gateway send/arrival time.\\
Step 4 & Verify qualified, unqualified and unavailable timing indications for their respective cases; check that unreliable or unknown time is never labeled qualified.\\
Step 5 & For delayed or repeated status, confirm that it still refers to the original sample time. Retain every mismatch and mapping dependency; do not assume a timestamp-accuracy bound not in this requirement.\\
Expected result & Every sample-related status preserves its source timestamp and correctly represents timing quality or unavailable time; no unqualified timestamp is represented as qualified. Execution is gated on the approved mission contract. Authoritative bound: \VerificationRef{VER-054}. This is an expectation, not an observed result.\\
Planned evidence & \ArtifactRef{Evidence}{EVD-054} (Not collected). Source samples and timing states, approved representation, correlated gateway captures, timestamp conversion comparison, unknown/unreliable timing cases and delayed/repeated-status results.\\
Record / restore & Record per-step observations and any deviation against the expected result, retain raw data and calculations, restore the baseline, and leave the result Pending until an executed review supplies a disposition.\\
Navigation & \hyperref[Table:WeekTwoTests]{Main procedure index}; \hyperref[Table:TraceabilityMatrix]{main trace matrix}.\\
\end{tabularx}
\normalsize
\subsection{Supporting Material}
\begin{landscape}
\small
\begin{tabularx}{\linewidth}{L{2cm}L{3cm}Y Y}
\caption{Supporting Main Baseline Planned Checks\label{Table:WeekTwoTestsSupporting}}\\
\toprule
\textbf{Check} & \textbf{Requirement / Use case} & \textbf{Procedure} & \textbf{Acceptance}\\
\midrule
\endfirsthead
\textbf{Check} & \textbf{Requirement / Use case} & \textbf{Procedure} & \textbf{Acceptance}\\
\midrule
\endhead
\midrule
\multicolumn{4}{r}{\small\itshape Continued on next page}\\
\endfoot
\bottomrule
\endlastfoot
\TestCaseID{TC-011} & \RequirementRef{REQ-F-004}\par \UseCaseRef{UC-01} & Run a known continuous source while processing and output are loaded; record read, processing, and delivery intervals. & Acceptance bound: \VerificationRef{VER-011}. Resolve all linked TBD dependencies first; otherwise no pass disposition.\\
\TestCaseID{TC-012} & \RequirementRef{REQ-F-005}\par \UseCaseRef{UC-01} & Vary conversion start, completion, and read timing; observe INTR, read control, and data bus with a logic analyzer. & Acceptance bound: \VerificationRef{VER-012}. Resolve all linked TBD dependencies first; otherwise no pass disposition.\\
\TestCaseID{TC-013} & \RequirementRef{REQ-F-006}\par \UseCaseRef{UC-02} & Inject incomplete conversions, incoherent capture, and unavailable timing reference; inspect record status. & Acceptance bound: \VerificationRef{VER-013}. Resolve all linked TBD dependencies first; otherwise no pass disposition.\\
\TestCaseID{TC-014} & \RequirementRef{REQ-F-007}\par \UseCaseRef{UC-01} & Apply reference voltages spanning the approved range; compare raw codes and reported voltage with the documented transfer function. & Acceptance bound: \VerificationRef{VER-014}. Resolve all linked TBD dependencies first; otherwise no pass disposition.\\
\TestCaseID{TC-015} & \RequirementRef{REQ-F-008}\par \UseCaseRef{UC-01} & Load approved and deliberately malformed, wrong-device, out-of-range, or expired calibration records. & Acceptance bound: \VerificationRef{VER-015}. Resolve all linked TBD dependencies first; otherwise no pass disposition.\\
\TestCaseID{TC-017} & \RequirementRef{REQ-I-002}\par \UseCaseRef{UC-01} & Inspect the versioned contract and decode representative records, including restart and sequence-boundary cases. & Acceptance bound: \VerificationRef{VER-017}. Resolve all linked TBD dependencies first; otherwise no pass disposition.\\
\TestCaseID{TC-018} & \RequirementRef{REQ-D-002}\par \UseCaseRef{UC-01} & Publish known streams through sequence exhaustion and restart, retaining epoch identifiers and records. & Acceptance bound: \VerificationRef{VER-018}. Resolve all linked TBD dependencies first; otherwise no pass disposition.\\
\TestCaseID{TC-020} & \RequirementRef{REQ-TIM-02}\par \UseCaseRef{UC-01} & Compare event timestamps with a calibrated independent clock under the approved load and synchronization conditions. & Acceptance bound: \VerificationRef{VER-020}. Resolve all linked TBD dependencies first; otherwise no pass disposition.\\
\TestCaseID{TC-021} & \RequirementRef{REQ-Q-002}\par \UseCaseRef{UC-01} & Acquire a stable input under the agreed load; independently measure consecutive conversion acceptance times. & Acceptance bound: \VerificationRef{VER-021}. Resolve all linked TBD dependencies first; otherwise no pass disposition.\\
\TestCaseID{TC-022} & \RequirementRef{REQ-Q-003}\par \UseCaseRef{UC-01} & Measure start-to-ready intervals across approved input, clock, supply, and temperature conditions. & Acceptance bound: \VerificationRef{VER-022}. Resolve all linked TBD dependencies first; otherwise no pass disposition.\\
\TestCaseID{TC-023} & \RequirementRef{REQ-Q-004}\par \UseCaseRef{UC-01} & Sweep independent calibration-check points, including endpoints and intermediate points not used to fit calibration. & Acceptance bound: \VerificationRef{VER-023}. Resolve all linked TBD dependencies first; otherwise no pass disposition.\\
\TestCaseID{TC-024} & \RequirementRef{REQ-Q-005}\par \UseCaseRef{UC-01} & Inject uniquely identifiable samples at the approved burst envelope from the specified initial queue state. & Acceptance bound: \VerificationRef{VER-024}. Resolve all linked TBD dependencies first; otherwise no pass disposition.\\
\TestCaseID{TC-025} & \RequirementRef{REQ-Q-006}\par \UseCaseRef{UC-01} & Analyze arrival and service envelopes and corroborate the bound using paused-consumer and burst tests. & Acceptance bound: \VerificationRef{VER-025}. Resolve all linked TBD dependencies first; otherwise no pass disposition.\\
\TestCaseID{TC-028} & \RequirementRef{REQ-Q-007}\par \UseCaseRef{UC-02} & Disconnect the selected sensor/USB path during acquisition and time the declared event to the observed status change. & Acceptance bound: \VerificationRef{VER-028}. Resolve all linked TBD dependencies first; otherwise no pass disposition.\\
\TestCaseID{TC-030} & \RequirementRef{REQ-Q-008}\par \UseCaseRef{UC-02} & Restore each approved recoverable fault and measure prerequisite-ready to first new qualified publication. & Acceptance bound: \VerificationRef{VER-030}. Resolve all linked TBD dependencies first; otherwise no pass disposition.\\
\TestCaseID{TC-031} & \RequirementRef{REQ-F-011}\par \UseCaseRef{UC-02} & Inject clock offset, drift, reset, and loss of synchronization; inspect timestamp and fault status. & Acceptance bound: \VerificationRef{VER-031}. Resolve all linked TBD dependencies first; otherwise no pass disposition.\\
\TestCaseID{TC-033} & \RequirementRef{REQ-F-012}\par \UseCaseRef{UC-04} & Drive a known count sequence through zero, maximum count, and wrap; compare captured words against an independent trace. & Acceptance bound: \VerificationRef{VER-033}. Resolve all linked TBD dependencies first; otherwise no pass disposition.\\
\TestCaseID{TC-034} & \RequirementRef{REQ-Q-009}\par \UseCaseRef{UC-04} & Sweep oscillator rate and host load while independently logging counter edges and capture times. & Acceptance bound: \VerificationRef{VER-034}. Resolve all linked TBD dependencies first; otherwise no pass disposition.\\
\TestCaseID{TC-035} & \RequirementRef{REQ-Q-010}\par \UseCaseRef{UC-04} & Compare captured-word timestamps against independently observed latch/capture events across the approved run duration. & Acceptance bound: \VerificationRef{VER-035}. Resolve all linked TBD dependencies first; otherwise no pass disposition.\\
\TestCaseID{TC-036} & \RequirementRef{REQ-I-003}\par \UseCaseRef{UC-04} & Capture around multi-bit transitions, cascaded carries, and wrap; observe latch control and all counter bits. & Acceptance bound: \VerificationRef{VER-036}. Resolve all linked TBD dependencies first; otherwise no pass disposition.\\
\TestCaseID{TC-037} & \RequirementRef{REQ-F-013}\par \UseCaseRef{UC-04} & Suppress captures and introduce gaps shorter and longer than a full counter cycle while retaining independent edge counts. & Acceptance bound: \VerificationRef{VER-037}. Resolve all linked TBD dependencies first; otherwise no pass disposition.\\
\TestCaseID{TC-038} & \RequirementRef{REQ-F-014}\par \UseCaseRef{UC-04} & Apply reset, enable/disable, known edge sequences, and wrap; compare the evaluation log with the independent trace. & Acceptance bound: \VerificationRef{VER-038}. Resolve all linked TBD dependencies first; otherwise no pass disposition.\\
\TestCaseID{TC-039} & \RequirementRef{REQ-Q-011}\par \UseCaseRef{UC-04} & Stop or perturb the oscillator and reference independently, including loss of the common comparison basis. & Acceptance bound: \VerificationRef{VER-039}. Resolve all linked TBD dependencies first; otherwise no pass disposition.\\
\TestCaseID{TC-040} & \RequirementRef{REQ-I-004}\par \UseCaseRef{UC-05} & Inspect the contract and parse frames containing every byte value, especially delimiter-like bytes and boundary lengths. & Acceptance bound: \VerificationRef{VER-040}. Resolve all linked TBD dependencies first; otherwise no pass disposition.\\
\TestCaseID{TC-041} & \RequirementRef{REQ-Q-012}\par \UseCaseRef{UC-05} & Transmit a known framed stream for the approved duration; measure payload bytes at the electrical and application boundaries. & Acceptance bound: \VerificationRef{VER-041}. Resolve all linked TBD dependencies first; otherwise no pass disposition.\\
\TestCaseID{TC-042} & \RequirementRef{REQ-F-015}\par \UseCaseRef{UC-05} & Force short writes, fragmented reads, timeouts, and cancellation while retaining submitted and accepted byte counts. & Acceptance bound: \VerificationRef{VER-042}. Resolve all linked TBD dependencies first; otherwise no pass disposition.\\
\TestCaseID{TC-043} & \RequirementRef{REQ-F-016}\par \UseCaseRef{UC-05} & Delete bytes and frames, duplicate frames, reorder frames, and reset the producer during known transfers. & Acceptance bound: \VerificationRef{VER-043}. Resolve all linked TBD dependencies first; otherwise no pass disposition.\\
\TestCaseID{TC-044} & \RequirementRef{REQ-Q-013}\par \UseCaseRef{UC-05} & Apply the approved producer burst while stalling and resuming the receiver; trace queue occupancy and end-to-end byte counts. & Acceptance bound: \VerificationRef{VER-044}. Resolve all linked TBD dependencies first; otherwise no pass disposition.\\
\TestCaseID{TC-045} & \RequirementRef{REQ-I-005}\par \UseCaseRef{UC-05} & Inspect wiring and directions; observe all bits and control signals through receiver pauses and alternating data patterns. & Acceptance bound: \VerificationRef{VER-045}. Resolve all linked TBD dependencies first; otherwise no pass disposition.\\
\TestCaseID{TC-046} & \RequirementRef{REQ-Q-014}\par \UseCaseRef{UC-05} & Measure correlated producer and consumer events with a common analyzer; perturb clock offset and drift. & Acceptance bound: \VerificationRef{VER-046}. Resolve all linked TBD dependencies first; otherwise no pass disposition.\\
\TestCaseID{TC-047} & \RequirementRef{REQ-D-003}\par \UseCaseRef{UC-05} & Inject defined single-bit and multi-bit corruptions in payload and protected headers and run approved check vectors. & Acceptance bound: \VerificationRef{VER-047}. Resolve all linked TBD dependencies first; otherwise no pass disposition.\\
\TestCaseID{TC-048} & \RequirementRef{REQ-F-017}\par \UseCaseRef{UC-05} & Disconnect mid-frame, corrupt length/framing, and restart either endpoint with queued data. & Acceptance bound: \VerificationRef{VER-048}. Resolve all linked TBD dependencies first; otherwise no pass disposition.\\
\TestCaseID{TC-049} & \RequirementRef{REQ-Q-015}\par \UseCaseRef{UC-05} & Measure accepted-frame to validated-delivery intervals using the approved common/correlated reference under normal and burst loads. & Acceptance bound: \VerificationRef{VER-049}. Resolve all linked TBD dependencies first; otherwise no pass disposition.\\
\end{tabularx}
\normalsize
\end{landscape}
\section{Retained Template Test Examples}
These service-request tests are template examples only.

\begin{tabularx}{\linewidth}{L{2.1cm}L{3.0cm}YL{5.0cm}}
\caption{Test Case Examples\label{Table:TestCaseExamples}}\\
\toprule
\textbf{ID} & \textbf{Mapped Requirement} & \textbf{Procedure} & \textbf{Expected Result}\\
\midrule
\endfirsthead
\toprule
\textbf{ID} & \textbf{Mapped Requirement} & \textbf{Procedure} & \textbf{Expected Result}\\
\midrule
\endhead
\TestCaseID{TC-001} & \RequirementRef{REQ-F-001} & Submit a request containing valid required data. & The system creates one request and displays a unique identifier.\\
\TestCaseID{TC-002} & \RequirementRef{REQ-F-002} & Open the status for a known request. & The current status and most recent update are displayed.\\
\TestCaseID{TC-003} & \RequirementRef{REQ-Q-001} & Measure 20 status requests under the reference workload. & Every measured response is displayed within two seconds.\\
\TestCaseID{TC-004} & \RequirementRef{REQ-F-003} & Export a known request set in the approved format. & The export contains the expected history without changing source records.\\
\TestCaseID{TC-005} & \RequirementRef{REQ-I-001} & Export representative request history and validate the result against the approved interface schema version. & The export conforms to the approved schema and declares the expected schema version.\\
\TestCaseID{TC-006} & \RequirementRef{REQ-S-001} & Attempt the same export as an authorized actor, an authenticated actor without the export role, and an unauthenticated actor. & Only the authorized actor receives an export; both unauthorized attempts are denied and recorded.\\
\TestCaseID{TC-007} & \RequirementRef{REQ-D-001} & Perform two status transitions and inspect the stored history and exported record. & Each transition retains the expected timestamp, actor identifier, previous status, and new status.\\
\TestCaseID{TC-008} & \RequirementRef{REQ-C-001} & Execute the critical use-case test suite on every platform and browser version in the approved deployment baseline. & Every critical use case passes on every approved platform and browser version.\\
\TestCaseID{TC-009} & \RequirementRef{REQ-B-001} & Time an authorized analyst preparing an approved export from the reference data set without manual reconstruction. & The analyst obtains analysis-ready request-history data within five minutes.\\
\bottomrule
\end{tabularx}
